\begingroup
\section{The established low-polynomial bound in the complete packet action}
\label{sec:LRP}

This section propagates the existing AKS26--32 and FEX31--33 estimates
through the exact inverse-dual map TAC6b and TAC9a. The estimate itself
belongs to those earlier calculations. The contribution here is its
explicit insertion into TAC10, TAC15--18 and FTR31, with all finite
constants, frames and four endpoint signs retained. The programme's
recorded origin and the creators of its deposited notes remain credited
\cite{human:globalization2026,human:splitzero2026}.

\subsection{The original metric and the finite bound}

Keep the actual hypothetical simple quartet, its fixed period and branch,
and the complete conductor of TAC1b. Write
\[
\begin{gathered}
 a=4l+1\ge9,\quad q=(a+1)^2,\quad s\in\{1,a\},\quad
 c=a/2,\quad R_a=a\sqrt{\delta^2+\gamma^2},\\
 \beta_{a,u,w}=c+(2u-a)\delta+i(2w-a)\gamma,\qquad
 \chi_a(S)=\prod_{u,w=0}^{a}(S-\beta_{a,u,w}),\\
 E_a=\mathbb C[S]/(\chi_a),\quad
 L_v=\mathbb C[S]_{<v}\subset E_a,\qquad 0\le v\le80.
\end{gathered}\tag{LRP1}
\]
The original inclusion of $L_v$ in the monomial coefficient frame of
$E_a$ is $P_{<v}$. Let $G_N^{(s)}$ be the quotient Gram of
$\mathbb C[S]_{\le N}\to E_a$, for the original Gamma source
\[
 dm_s(y)=\frac{M_s2^{s/2}}{4\pi\Gamma(s/2)}
    |\Gamma(s/4+iy/2)|^2\,dy,\qquad M_s=(2\pi)^{s/2},\quad S=c+iy.
 \tag{LRP2}
\]
Thus $G_{L,N}^{(s)}=P_{<v}^*G_N^{(s)}P_{<v}$. The map and metric are
the ones in FEX31 and TAC6b; the target's centre $c-4$ is not substituted
for the centre $c$ of this whole numerator packet.

For completeness the finite covariance estimate needed here is derived
on the actual source. Set $b_s=s/2$. Its original orthogonal polynomials
and their norms are
\[
\begin{gathered}
 p_0=1,\quad p_1=y,\quad
 p_{d+1}=yp_d-d(d-1+b_s)p_{d-1},\qquad h_{d,s}=M_s d!(b_s)_d,\\
 u_d(S)=p_d((S-c)/i)/\sqrt{h_{d,s}},\qquad
 \sum_{d\ge0}\frac{p_d(y)}{d!}z^d
       =(1+z^2)^{-b_s/2}e^{y\arctan z}.
\end{gathered}\tag{LRP3}
\]
The precise classical dictionary is
$p_d(y)=d!P_d^{(s/4)}(y/2;\pi/2)$.
The authored NIST DLMF formulas give its weight, norms and generating
function, including this change of variable
\cite[eqs.~(18.19.8)--(18.19.9), (18.23.7)]{human:dlmf18}.
The mass and phase calculation is given in HPS and TAC1a; AKS15--16
also derives the generating function from this recurrence and the norm
from the original transform
$\int e^{ty}dm_s=M_s(\cos t)^{-b_s}$, $|t|<\pi/2$.

Let $C_s$ have the original monomial coefficient columns
$[u_0],\ldots,[u_{q-1}]$. It is invertible, with determinant
$i^{-q(q-1)/2}\prod_{d<q}h_{d,s}^{-1/2}$. Put
\[
 v_j=C_s^{-1}[u_{q+j}],\qquad
 \mathsf K_j=I_q+\sum_{r=0}^{j-1}v_rv_r^*,\qquad
 T_s=C_s^{-1}P_{<v},\qquad 0\le j\le q+1.
 \tag{LRP4}
\]
The full source map in its orthonormal basis is
$C_s[I_q,v_0,\ldots,v_{j-1}]$. Its minimum lift of $C_sx$ is
$[I_q,v_0,\ldots,v_{j-1}]^*\mathsf K_j^{-1}x$:
it has the required image and is orthogonal to the full relation kernel.
Pythagoras therefore proves
\[
 G_{q+j-1}^{(s)}=C_s^{-*}\mathsf K_j^{-1}C_s^{-1},\qquad
 H_j:=G_{L,q+j-1}^{(s)}=T_s^*\mathsf K_j^{-1}T_s.
 \tag{LRP5}
\]
This is the original AKS26 compression, with the literal inclusion
$P_{<v}$ in place of a general fixed kernel frame.

Here are the complete finite constants from AKS3 and AKS30, with
$\kappa_0$ denoting AKS30's scalar $c_0$, to avoid confusion with the
source centre:
\[
\begin{gathered}
 t_q=\pi/2-1/q,\quad \alpha_q=2\log3/t_q,\quad
 \Sigma_q=(\alpha_q^q-1)/(\alpha_q-1),\\
 B_{s,q}=[(9/25)\sin(1/q)]^{-s/2}
           \exp\{2R_a(\log3+t_q/2)\},\\
 \mathcal L_{s,j}=
 \log\left(1+B_{s,q}\Sigma_q^2
        \sum_{d=q}^{q+j-1}\frac{d!}{(s/2)_d}(25/16)^d\right),\\
 \alpha_\infty=4\log3/\pi,\qquad
 \kappa_0=\log\frac{25\log3}{4\pi}>0,\\
 E_{s,q}(R_a)=\log2+\frac4{\pi-2/q}
       +\frac{s}{2}\log\frac{25\pi q}{18}
       +2R_a\left(\log3+\frac\pi4\right)\\
 \hspace{21mm}+\log((q+1)(4q+1))-2\log(\alpha_\infty-1).
\end{gathered}\tag{LRP6}
\]
An empty sum has value zero. To verify the covariance bound including
every root coefficient, put $P(y)=i^{-q}\chi_a(c+iy)$ and list its
roots $\omega_\nu=(\beta_\nu-c)/i$ in their original order. All satisfy
$|\omega_\nu|\le R_a$. Newton interpolation gives, for complex $z$,
\[
 \operatorname{rem}_P e^{zy}=
 \sum_{r=0}^{q-1}[\omega_1,\ldots,\omega_{r+1}]e^{z\cdot}
       \prod_{\nu=1}^{r}(y-\omega_\nu).
\]
The divided difference equals the integral of
$z^r e^{z\sum\theta_\nu\omega_\nu}$ over the standard $r$-simplex
of volume $1/r!$. This identity follows by expanding the exponential:
the simplex moments are $\prod n_\nu!/(r+\sum n_\nu)!$, and the
Newton divided difference of $y^{r+n}$ is the complete homogeneous
polynomial of degree $n$ in the nodes. Hence its absolute value is at
most $|z|^r e^{R_a|z|}/r!$.

For $0<t<\pi/2$ the positive cosh series in the original transform gives
$\|y^r\|\le\sqrt{M_s}(\cos t)^{-s/4}(2/t)^r r!$.
Expanding the entire Newton product and bounding each elementary
symmetric coefficient gives
\[
 \left\|\prod_{\nu=1}^{r}(y-\omega_\nu)\right\|
 \le\sqrt{M_s}(\cos t)^{-s/4}(2/t)^r r!e^{R_at/2}.
\]
For a complex circle $|w|=z<1$, use
$|\arctan w|\le\operatorname{arctanh}z$ and
$|1+w^2|\ge1-z^2$ in the finite remainder of LRP3.
Cauchy coefficient extraction, followed by the preceding full-vector
norm bound and division by $h_{d,s}$, yields
\[
\begin{split}
 \|v_{d-q}\|^2\le{}&\frac{d!}{(s/2)_d}z^{-2d}
 [(1-z^2)\cos t]^{-s/2}
 e^{2R_a(\operatorname{arctanh}z+t/2)}\\
 &\mathrel{}\times
 \left[\sum_{r=0}^{q-1}
       (2\operatorname{arctanh}z/t)^r\right]^2.
\end{split}\tag{LRP7}
\]
This extraction is performed on a finite remainder vector holomorphic
in $|w|<1$; it requires no integral of the generating function over
the original source on that large circle. Parseval in its low-degree
orthonormal basis identifies the left side with the full remainder
norm divided by $h_{d,s}$. The two original masses thus cancel by
this exact equality. Setting $z=4/5$, $t=t_q$ and summing the complete
positive rank-one matrices proves
\[
 I_q\preceq\mathsf K_j\preceq e^{\mathcal L_{s,j}}I_q.
 \tag{LRP8}
\]
The finite estimate $\mathcal L_{s,j}\le2\kappa_0q+E_{s,q}(R_a)$
for $j=q,q+1$ follows as in AKS31. Explicitly,
$d!/(s/2)_d\le4^d/\binom{2d}{d}\le2d+1$, so the sum in LRP6
is at most $(q+1)(4q+1)(5/4)^{4q}$. Also
\[
 \Sigma_q\le\frac{\alpha_q^q}{\alpha_\infty-1},\quad
 2q\log\alpha_q\le2q\log\alpha_\infty+\frac4{\pi-2/q},\quad
 \sin(1/q)\ge\frac2{\pi q}.
\]
The first holds since $\alpha_q\ge\alpha_\infty>1$; the second
uses $-\log(1-x)\le x/(1-x)$, and the third is the sine chord bound.
The exponential in LRP6 uses $t_q\le\pi/2$. The order-$q$ terms add
to $2\kappa_0q$, and $\log(1+e^x)\le\log2+\max(0,x)$ gives the
stated finite bound. Here $q\ge100$ satisfies AKS31's guard $q\ge36$.

For $v>0$ set $d_j=-\log\det(H_0^{-1/2}H_jH_0^{-1/2})$.
Inversion and compression in LRP8 show
$0=d_0\le d_1\le\cdots\le d_{q+1}$ and
$d_j\le v\mathcal L_{s,j}$. Define
\[
\begin{gathered}
 V_s=\|v_0\|^2,\qquad
 \theta_s=v_0^*T_s(T_s^*T_s)^{-1}T_s^*v_0,\\
 \ell_{a,s}=\log\frac{1+V_s}{1+V_s-\theta_s}=d_1,\qquad
 \mathcal B^{L}_{a,s}=v[4\kappa_0q+2E_{s,q}(R_a)],\qquad
 u_{a,s}=\mathcal B^{L}_{a,s}-\ell_{a,s}.
\end{gathered}\tag{LRP9}
\]
The matrix defining $\theta_s$ is the orthogonal projection onto
$\operatorname{im}T_s$, so $0\le\theta_s\le V_s$.
The inverse identity
$(I+v_0v_0^*)^{-1}=I-v_0v_0^*/(1+V_s)$ and its compressed
rank-one determinant prove the formula for $d_1$.
For $v=0$ set $\ell_{a,s}=u_{a,s}=\mathcal B^L_{a,s}=0$;
all low determinants then equal one.

The exact primal four-endpoint return is
\[
\begin{split}
 F_{L_a}^{(s)}
 &=\log\frac{\det G_{L,q-1}^{(s)}\det G_{L,q}^{(s)}}
                  {\det G_{L,2q-1}^{(s)}\det G_{L,2q}^{(s)}}
   =d_q+d_{q+1}-d_1,\\
 &\hspace{7mm}\boxed{\quad
 0\le\ell_{a,s}\le F_{L_a}^{(s)}\le u_{a,s}
       \le\mathcal B^L_{a,s}.\quad}
\end{split}\tag{LRP10}
\]
The lower bound uses $d_q,d_{q+1}\ge d_1$ and the upper bound uses
LRP8 and LRP6 at both high increments. This retains the sharpened
AKS28 interval in addition to FEX32's bound.

For fixed quartet and conductor, $R_a=O_{\delta,\gamma}(a)$ and
$v\le80$. The displayed constants give
$E_{1,q}(R_a)=O_{\delta,\gamma}(a+\log q)$ and
$E_{a,q}(R_a)=O_{\delta,\gamma}(a\log q)$. Consequently
\[
 0\le F_{L_a}^{(s)}\le\mathcal B^L_{a,s}
       =O_{\delta,\gamma}(q)=o(aq),\qquad s=1,a.
 \tag{LRP11}
\]
This holds along the original discrete family; its finite value is
never set to zero. The difference between source orders has the
directed bound
\[
 \ell_{a,a}-u_{a,1}
 \le F_{L_a}^{(a)}-F_{L_a}^{(1)}
 \le u_{a,a}-\ell_{a,1}.
 \tag{LRP12}
\]
No sign for that difference is required or inferred.

\subsection{The exact dual map and propagation through the receiving equations}

TAC6b proves on the complete original dual packet that
$Q_{v,N}^{(s)}=\overline{(G_{L,N}^{(s)})^{-1}}$ in the paired
coefficient frames. Its proof minimizes the full dual norm over the
annihilator of $L_v$; this is the inverse restriction metric.
The Schur minimum and block factorization used there agree with
the classical positive-matrix formulas
\cite[Appendices A.5.5, p.~650, and C.4.1, p.~673]{human:boyd-vandenberghe}.
Since each determinant is positive real, its logarithm changes sign
under this conjugate inverse. With the original dual sign operator
$\mathcal R f=-f(q-1)-f(q)+f(2q-1)+f(2q)$ this proves
\[
 \mathcal V_{a,s}=\mathcal R\log\det Q_{v,N}^{(s)}
       =-\mathcal R\log\det G_{L,N}^{(s)}=F_{L_a}^{(s)}.
 \tag{LRP13}
\]
Thus LRP10 bounds the actual final factor in TAC10, without changing
the analytic middle quotient or the conductor restriction.

Retain the exact complete action TAC15 and put
\[
 A^0_{a,s,R}=\mathcal C_{a,s}+\delta_{a,s}^{\rm ar}
            +\mathscr T_{a,s,R}+\mathcal W_q^\sigma.
 \tag{LRP14}
\]
All four terms keep TAC's original metric and cutoffs. In particular
$\delta_{a,1}^{\rm ar}=\delta_a^\sigma$ and
$\delta_{a,a}^{\rm ar}=\mathcal H_a+\delta_a^\sigma$.
Substituting LRP10 in the two sides of TAC16 gives
\[
\begin{split}
 A^0_{a,s,R}+\ell_{a,s}+\Xi^-_{a,s}-E_{\rm hi}-\Pi_{h,a}
       &\le J_a,\\
 J_a&\le A^0_{a,s,R}+u_{a,s}+\Xi^+_{a,s}
                   +\mathfrak E_{s,R}+E_{\rm lo}.
\end{split}\tag{LRP15}
\]
This is a finite interval for the same action, proved by scalar
inequalities on the same actual matrices. At $a=j=k+4$ and the
original $R_j$, FTR28's three joint controls give the further upper bound
\[
\begin{split}
 J_j\le{}&\mathcal C_{j,s}+\delta_{j,s}^{\rm ar}
          +\mathcal W_{q_j}^{\sigma}+\mathfrak U_R+u_{j,s}\\
         &+\Xi^+_{j,s}+\mathfrak E_{s,R}+E_{\rm lo}.
\end{split}\tag{LRP16}
\]
Indeed substitute $\mathscr T_{j,s,R}\le\mathfrak U_R$ in LRP15.
Every coefficient remains one, exactly as in FTR31.
Since $jq_j/(kq_k)\to1$, LRP11 also controls this term on the
original comparison scale $kq_k$.

For the allocation identity TAC17 set
$\mathcal E_{a,s}^{\rm src}=\sum_{i=0}^3\sigma_i e_{N_i,s}$,
with $\sigma=(1,1,-1,-1)$, and keep the actual observation kernel
$K_a$ and $t_0=\dim(K_a\cap L_v)$. FEX32 passes the same full
covariance inequality by restriction and minimization to
$(K_a+L_v)/K_a$, of dimension $v-t_0$. Therefore
\[
\begin{split}
 -u_{a,s}\ \le\
 \delta_{K,a,s}-\mathcal E_{a,s}^{\rm src}
   &=-F_{L_a}^{(s)}+F_{(K_a+L_v)/K_a}^{(s)}\\
   &\le(v-t_0)(\mathcal L_{s,q}+\mathcal L_{s,q+1})-\ell_{a,s}.
\end{split}\tag{LRP17}
\]
Both bounds are $O_{\delta,\gamma}(q)$. All source Jacobians remain
in $\mathcal E_{a,s}^{\rm src}$; this statement makes no estimate
for them. At $s=a$, TAC18 still cancels the exact two occurrences of
$\delta_{K,a,a}$ and the exact five allocation terms with coefficients
$+1/4$ and $-1/4$. Applying separate upper bounds before that identity
is unnecessary; the cancellation remains exact.

The two original Gamma source orders also give a useful exact check
on the same arithmetic action. Write $\Delta X=X_{a,a}-X_{a,1}$
for its source-order difference and retain the whole native return
$\mathcal T_{a,s}=\mathcal T_{a,s}^{\le R}+\mathscr T_{a,s,R}$.
Subtracting the two TAC15 identities cancels $J_a$, $W_a$ and both
arithmetic penalties. TAC13 gives
$\delta_{a,a}^{\rm ar}-\delta_{a,1}^{\rm ar}=\mathcal H_a$.
Hence
\[
\begin{gathered}
 \Delta\mathcal C+\Delta\mathcal T+\Delta\Xi+\mathcal H_a
       =-\bigl(F_{L_a}^{(a)}-F_{L_a}^{(1)}\bigr),\\
 \ell_{a,1}-u_{a,a}
 \le\Delta\mathcal C+\Delta\mathcal T+\Delta\Xi+\mathcal H_a
 \le u_{a,1}-\ell_{a,a}.
\end{gathered}\tag{LRP18}
\]
The combined expression is $O_{\delta,\gamma}(q)=o(aq)$ by LRP11.
Its conductor, native and source-correction terms retain their exact
values and correlation. The calculation compares source orders at the
same packet degree $a$; it does not substitute for the separate
degree-$k$ proper-source return of PRD and NFR.

The cumulative receiving calculation therefore has an evaluated
subleading bound for its low-polynomial factor. Its exact value,
the independent proper degree-$k$ source, the complete conductor
restriction $\mathcal C_{a,s}$, the correlated native tail, the
arithmetic correction and all nine action blocks remain in their
proved formulas. LRP11 evaluates only this identified low factor;
the finite controls for the remaining factors retain their actual values.
\endgroup
